\frontsection{Conclusion Générale}

Ce stage a permis de concevoir et de développer une plateforme web complète
de gestion associative, couvrant l'ensemble des besoins exprimés dans le
cahier des charges : gestion des adhérents et de leurs cotisations annuelles,
suivi d'un solde dû par adhérent et par année, gestion des projets
associatifs à travers un cycle de vie complet (de la création à la clôture
ou à l'annulation), constitution de comités dédiés avec des droits de saisie
propres à chaque rôle, workflow d'approbation des dons et des dépenses avec
justificatifs, génération automatique des rapports de clôture de projet et
des rapports association, journal d'activité et centre de notifications, et
consultation en libre-service par chaque adhérent de son propre historique.

Le système repose sur une architecture web en trois niveaux (interface React,
API REST Laravel, base de données MySQL), avec une autorisation à deux axes
— rôle global dans l'association et rôle au sein d'un comité de projet — qui
reflète fidèlement la manière dont une association de ce type fonctionne
réellement : un même adhérent peut être simple cotisant sur l'ensemble de
l'association et trésorier d'un comité sur un projet particulier. Cette
séparation, au cœur du cahier des charges, structure l'ensemble du modèle de
données et du système de permissions présentés au chapitre~3.

Ce projet a permis de consolider des compétences techniques en conception
d'API REST sécurisées, en modélisation de workflows métier (cycles de vie,
approbations à plusieurs états) et en développement d'interfaces web
réactives, ainsi que des compétences plus larges : traduire un besoin
métier exprimé en langage naturel en un modèle de données et un système de
permissions cohérents, et itérer sur ce modèle à mesure que des cas
particuliers apparaissent.

Plusieurs pistes d'évolution restent ouvertes pour la suite du projet :
une application mobile pour les adhérents, afin de faciliter la consultation
de leur historique et le paiement de leur cotisation ; des notifications par
SMS pour les adhérents ne disposant pas d'un smartphone ou d'un accès
régulier à internet, en complément des notifications actuelles ; et un
support multi-association, pour permettre à la même plateforme de servir
plusieurs associations indépendantes plutôt qu'une seule instance dédiée.
